% Hodina 41: Algebraický výraz a mnohočlen. Oprava 2026-10-30: jednočlen a cesky zapis intervalu.
% Výchozí příprava: Werdna09/Work-MAT, 1_rocnik/pripravy/41_ALG_PojmyAMnohocleny.tex.
% Učebnice: ucebnice/01_Zakladni_poznatky/Algebraicke_vyrazy/Algebraické výrazy.tex.
% Metodická inspirace: Realisticky.cz, Výrazy – definiční obor; Sčítání a odčítání mnohočlenů.
% Navazuje hodina 42: Sčítání a odčítání mnohočlenů.
\documentclass[11pt,a4paper]{article}
% Spolecny styl pro ucitelske pripravy MAT (Work/MAT)
% Verze 1.1, 2026-10-23
\NeedsTeXFormat{LaTeX2e}
\ProvidesPackage{priprava}[2026/10/23 v1.1 Sablona ucitelskych priprav MAT]

% Zakladni typografie a cestina
\RequirePackage[T1]{fontenc}
\RequirePackage[utf8]{inputenc}
\RequirePackage[czech]{babel}
\RequirePackage{lmodern}
\RequirePackage{microtype}
\RequirePackage{amsmath,amssymb}
\RequirePackage{xcolor}
\RequirePackage{enumitem}
\RequirePackage[most]{tcolorbox}
\RequirePackage{tabularx,booktabs,array}
\RequirePackage{fancyhdr}
\RequirePackage{needspace}
\RequirePackage{geometry}
\RequirePackage[hidelinks]{hyperref}

\geometry{a4paper,top=17mm,bottom=18mm,left=23mm,right=23mm}
\definecolor{notegray}{gray}{0.39}
\setlength{\parindent}{0pt}
\setlength{\parskip}{2.5pt plus 1pt minus 0.5pt}
\setlist[itemize]{topsep=2pt,itemsep=1pt,parsep=0pt,leftmargin=1.45em}
\setlist[enumerate]{topsep=2pt,itemsep=2pt,parsep=0pt,leftmargin=1.65em}
\setcounter{secnumdepth}{0}
\renewcommand{\arraystretch}{1.12}

% Zahlavi: vlevo rocnik | globalni tema, vpravo cislo: nazev hodiny
% Zapati zustava zcela prazdne.
\pagestyle{fancy}
\fancyhf{}
\renewcommand{\headrulewidth}{0.35pt}
\renewcommand{\footrulewidth}{0pt}
\setlength{\headheight}{14pt}
\newcommand{\nastavPripravu}[4]{%
  \fancyhead[L]{\small #1\enspace\textbar\enspace #2}%
  \fancyhead[R]{\small #3: #4}%
}

% Minimalisticke cernobile boxy. Kazda definice zvlast.
\tcbset{commonlesson/.style={enhanced,breakable=false,colback=white,colframe=black,
  colbacktitle=white,coltitle=black,boxrule=0.45pt,
  arc=0.8mm,outer arc=0.8mm,left=3mm,right=3mm,top=1.6mm,bottom=1.6mm,
  toptitle=1mm,bottomtitle=0.7mm,fonttitle=\bfseries,
  attach boxed title to top left={xshift=2.5mm,yshift=-1.15mm},
  boxed title style={colback=white,colframe=white,boxrule=0pt,left=1mm,right=1mm}}}
\newtcolorbox{zadaniBox}[1][]{commonlesson,title={Zadání},#1}
\newtcolorbox{prikladBox}[1][]{commonlesson,title={Příklad},#1}
\newtcolorbox{definiceBox}[1][]{commonlesson,title={Definice},#1}

% Spolecny box pro zadani a jeho reseni, bez mezery mezi castmi.
% Pouziti:
%   \begin{prikladReseniBox}
%     Zadani ulohy.
%     \tcblower
%     \begin{align*} ... \end{align*}
%   \end{prikladReseniBox}
% Spodni cast ma jemny sedy podklad a delici linku; lze v ni sazet
% vice radku matematiky. Obe casti maji jen jeden spolecny ramecek.
\tcbset{solutionlesson/.style={
  commonlesson,
  enhanced jigsaw,
  breakable,
  colbacklower=black!3,
  segmentation style={solid,draw=black!25,line width=0.35pt},
  fontlower=\small\color{notegray},
  before lower={\noindent\textbf{Řešení:}\enspace},
  leftlower=3mm,rightlower=3mm
}}
\newtcolorbox{prikladReseniBox}[1][]{solutionlesson,title={Příklad},#1}
\newtcolorbox{zadaniReseniBox}[1][]{solutionlesson,title={Zadání},#1}

% Pomucky pro autora pripravy: sedive poznamky; \reseni pro starsi pripravy
\newcommand{\poznamka}[1]{\par\begingroup\small\color{notegray}\itshape
  \textnormal{\textbf{Poznámka pro učitele.}}\enspace #1\par\endgroup}
\newcommand{\reseni}[1]{\par\begingroup\small\color{notegray}
  \textbf{Očekávané řešení:}\enspace #1\par\endgroup}
\newcommand{\blok}[1]{\needspace{5\baselineskip}\subsection*{#1}}
\newcommand{\vyrok}[1]{\uv{#1}}

\nastavPripravu{1. ročník}{Algebraické výrazy}{41}{Algebraický výraz a mnohočlen}
\begin{document}

\blok{Výklad a souvislosti}

Výrazy $3x+5$, $2a^2b$ nebo $\frac1x$ již umíme vyčíslit. Nyní přesněji pojmenujeme jejich části, rozlišíme jednočleny a mnohočleny a naučíme se určovat stupeň mnohočlenu.

\begin{prikladReseniBox}[title={Úvodní společný příklad},breakable=false]
Pro $P(x)=3x^2-2x+5$ vypočítej $P(2)$. Která čísla násobí jednotlivé mocniny $x$? Který člen neobsahuje proměnnou?
\tcblower
$P(2)=3\cdot 2^2-2\cdot 2+5=13$. Členy jsou $3x^2$, $-2x$, $5$; jejich koeficienty jsou $3$, $-2$, $5$. Nejvyšší mocnina s nenulovým koeficientem je $x^2$.
\end{prikladReseniBox}

\poznamka{Na stejném výrazu postupně zaváděj pojmy. Znaménko minus patří ke členu $-2x$; výraz lze chápat jako součet $3x^2+(-2x)+5$.}

\blok{Základní pojmy a definice}

\begin{definiceBox}[title={Definice: Algebraický výraz}]
\textbf{Algebraický výraz} je smysluplný matematický zápis vytvořený z čísel a proměnných pomocí konečného počtu algebraických operací (sčítání, odčítání, násobení, dělení, umocňování a odmocňování), případně závorek.
\end{definiceBox}

\begin{definiceBox}[title={Definice: Proměnná}]
\textbf{Proměnná} je symbol, kterému lze přiřazovat hodnoty z určené číselné množiny; například $x\in\mathbb R$.
\end{definiceBox}

\begin{definiceBox}[title={Definice: Definiční obor výrazu}]
\textbf{Definiční obor výrazu} $V(x)$ je množina všech hodnot proměnné, pro které mají všechny operace ve výrazu smysl v uvažovaném číselném oboru. Značíme jej $D_V$.
\end{definiceBox}

Například pro $V(x)=\frac1x$ je $D_V=\mathbb R\setminus\{0\}$, zatímco pro $W(x)=\sqrt{x}$ je při práci v reálných číslech $D_W=\langle 0;+\infty)$.

\begin{definiceBox}[title={Definice: Jednočlen}]
\textbf{Jednočlen (monom)} je výraz tvořený součinem reálného čísla a mocnin proměnných s nezápornými celočíselnými exponenty. Například pro dvě proměnné má tvar
\[
M=ax^m y^n,\qquad a\in\mathbb R,\quad m,n\in\mathbb N_0.
\]
Jednočlen neobsahuje součet ani rozdíl různých členů.
\end{definiceBox}

\begin{prikladReseniBox}[title={Příklad: Jednočlen, nebo mnohočlen?},breakable=false]
Roztřiď výrazy $-ab$, $3x$, $\sqrt7$, $-5a^2b^3$, $3x+2$ a $x^2-1$.
\tcblower
Jednočleny: $-ab$, $3x$, $\sqrt7$, $-5a^2b^3$ (také konstanty jsou jednočleny). Výrazy $3x+2$ a $x^2-1$ obsahují součet či rozdíl více jednočlenů, jsou to mnohočleny, ale nikoli jednočleny.
\end{prikladReseniBox}

\begin{definiceBox}[title={Definice: Koeficient jednočlenu}]
\textbf{Koeficient jednočlenu} je číselný činitel $a$ v součinu $M=ax^m y^n$. Například u $-ab$ je koeficient $-1$, u $3x$ je $3$.
\end{definiceBox}

\begin{definiceBox}[title={Definice: Stupeň jednočlenu}]
\textbf{Stupeň nenulového jednočlenu} je součet exponentů všech jeho proměnných. Pro $M=ax^m y^n$, $a\ne0$, platí
\[
\deg M=m+n.
\]
Nenulová konstanta má stupeň $0$.
\end{definiceBox}

\poznamka{Vysvětli, že $ax^m y^n$ znamená násobení, nikoli sčítání. Například $-5a^2b^3$ je jediný jednočlen stupně $5$. Součet $3x+2$ už je dvojčlen. Stupeň nulového jednočlenu neurčujeme.}

\begin{definiceBox}[title={Definice: Mnohočlen}]
\textbf{Mnohočlen (polynom)} je konečný součet jednočlenů. Každý jednočlen je tedy také mnohočlenem, ale mnohočlen s více nenulovými členy není jednočlenem.
\end{definiceBox}

\begin{definiceBox}[title={Definice: Obecný tvar mnohočlenu jedné proměnné}]
Každý \textbf{nenulový mnohočlen} v jedné proměnné $x$ lze jednoznačně zapsat ve tvaru
\[
P(x)=a_0+a_1x+a_2x^2+\dots+a_nx^n,
\]
kde $a_0,a_1,\ldots,a_n\in\mathbb R$, $a_n\ne0$ a $n\in\mathbb N_0$.
\end{definiceBox}

\begin{definiceBox}[title={Definice: Člen mnohočlenu}]
\textbf{Členem mnohočlenu} v zápisu $P(x)=\sum_{i=0}^{n}a_ix^i$ je jednočlen $a_ix^i$; číslo $a_i$ je \textbf{koeficient při mocnině} $x^i$. Číslo $a_0$ nazýváme \textbf{konstantním členem}.
\end{definiceBox}

\begin{definiceBox}[title={Definice: Stupeň mnohočlenu}]
\textbf{Stupeň nenulového mnohočlenu} $P(x)=a_0+a_1x+\cdots+a_nx^n$ s $a_n\ne0$ je číslo $n$; píšeme $\deg P=n$. U mnohočlenu s více proměnnými je stupeň největší stupeň jeho nenulového jednočlenu po sloučení stejných členů.
\end{definiceBox}

Nulový mnohočlen $P(x)=0$ nemá v této přípravě definovaný stupeň. U víceproměnného mnohočlenu například $2x^3y^2-xy+1$ je stupeň $5$, protože člen $2x^3y^2$ má stupeň $3+2$.

\poznamka{Zdůrazni rozdíl: $\sqrt2\,x^2$ je jednočlen, a tedy i mnohočlen (odmocnina je jen číslo), ale $\sqrt x$ a $1/x$ nejsou mnohočleny. Obecný tvar s $a_n\ne0$ popisuje pouze nenulové mnohočleny.}

\blok{Řešené příklady}

\begin{prikladReseniBox}[title={Příklad 1},breakable=false]
Urči členy, koeficienty a stupeň mnohočlenu $P(x)=3x^4-2x^2+7$.
\tcblower
Členy: $3x^4$, $-2x^2$, $7$; nenulové koeficienty: $a_4=3$, $a_2=-2$, $a_0=7$. Chybějící koeficienty jsou $a_3=a_1=0$. Platí $\deg P=4$.
\end{prikladReseniBox}

\begin{prikladReseniBox}[title={Příklad 2},breakable=false]
Urči stupeň mnohočlenu $Q(x,y)=2x^3y^2-xy+1$.
\tcblower
Stupně členů jsou $3+2=5$, $1+1=2$ a $0$. Proto $\deg Q=5$.
\end{prikladReseniBox}

\blok{Samostatné procvičení}

\begin{zadaniReseniBox}[title={Procvičení},breakable=false]
\begin{enumerate}[label=\alph*)]
\item Rozhodni, zda $x^2+\frac1x$ a $x^2-\frac34x+\sqrt3$ jsou mnohočleny.
\item Urči koeficient a stupeň jednočlenu $-5a^2b^3$.
\item Urči stupeň mnohočlenu $7$ a vysvětli případ $0$.
\end{enumerate}
\tcblower
\begin{enumerate}[label=\alph*)]
\item První není mnohočlen (obsahuje $x^{-1}$); druhý je mnohočlen druhého stupně.
\item Koeficient $-5$, stupeň $2+3=5$.
\item $\deg 7=0$; stupeň nulového mnohočlenu zde není definován.
\end{enumerate}
\end{zadaniReseniBox}

\blok{Navazující procvičování ve třídě}

\begin{zadaniReseniBox}[title={Základ: rozpoznání a vyčíslení},breakable=false]
\begin{enumerate}[label=\alph*)]
\item Pro $P(x)=-x^3+\frac12x^2-4$ urči členy, koeficienty a stupeň; vypočítej $P(0)$ a $P(-2)$.
\item Urči koeficient a stupeň $3a^2b^4$, $-xy^3$, $\frac12x$ a $-8$.
\item Rozhodni, které výrazy jsou mnohočleny: $\frac{x^2}{3}+1$, $\frac3x+1$, $\sqrt5x^2-x$, $\sqrt x+1$.
\end{enumerate}
\tcblower
\begin{enumerate}[label=\alph*)]
\item Členy $-x^3$, $\frac12x^2$, $-4$; koeficienty $-1$, $\frac12$, $-4$; stupeň $3$. $P(0)=-4$, $P(-2)=6$.
\item Po řadě $(3,6)$, $(-1,4)$, $(\frac12,1)$, $(-8,0)$.
\item Ano, ne, ano, ne.
\end{enumerate}
\end{zadaniReseniBox}

\begin{zadaniReseniBox}[title={Rezerva: zápis může klamat},breakable=false]
\begin{enumerate}[label=\alph*)]
\item Urči stupeň $3x^4-3x^4+2$ a $2x^2y^3+xy^5$.
\item Obdélník má délky stran $x+2$ a $3$ centimetry, $x>0$. Zapiš jeho obsah a vypočítej jej pro $x=4$.
\end{enumerate}
\tcblower
\begin{enumerate}[label=\alph*)]
\item První výraz je konstanta $2$, tedy stupeň $0$; druhý má stupeň $6$.
\item $S=3(x+2)\,\mathrm{cm}^2$; pro $x=4$ je $S=18\,\mathrm{cm}^2$.
\end{enumerate}
\end{zadaniReseniBox}

\poznamka{Ve 45minutové hodině jsou prioritou definice s příklady a základní procvičení; rezervu lze vynechat. Systematické slučování členů patří do hodiny 42.}

\blok{Domácí procvičení}

\begin{zadaniReseniBox}[title={Domácí úlohy},breakable=false]
\begin{enumerate}[label=\alph*)]
\item U $Q(x)=2x^3-x+5$ urči členy, koeficienty, stupeň a $Q(-1)$.
\item Urči stupeň $-4a^3b^2+2ab^4+1$.
\item Rozhodni, zda $x/5+2$, $5/x+2$ a $\sqrt2x+1$ jsou mnohočleny.
\item Urči stupeň konstantních mnohočlenů $9$ a $0$ podle naší konvence.
\end{enumerate}
\tcblower
\begin{enumerate}[label=\alph*)]
\item Členy $2x^3$, $-x$, $5$; koeficienty $2$, $-1$, $5$; stupeň $3$; $Q(-1)=4$.
\item Stupeň $5$.
\item Ano, ne, ano.
\item Stupeň $9$ je $0$; stupeň nulového mnohočlenu není definován.
\end{enumerate}
\end{zadaniReseniBox}

\end{document}
